\section{WorldEpisode Architecture}
\label{sec:architecture}

WorldEpisode is implemented as a sidecar blueprint around existing formats. It does not ask a
LeRobot dataset to become a USD scene, or a USD scene to become a policy-training table. Instead,
it records the relationships those containers need in order to agree about the same episode: the
world revision, persistent entities, space/time conventions, action contract, asset roles,
interaction events, provenance, and split lineage.

Figure~\ref{fig:worldepisode-bridge} shows the intended data flow. Dataset and log containers plug
in on the left; world assets and replay targets plug in on the right; the WorldEpisode sidecar
holds the semantics that would otherwise be scattered across undocumented loader assumptions.

\begin{figure}[t]
\centering
\resizebox{\linewidth}{!}{%
\begin{tikzpicture}[font=\sffamily]
  \node[wecard=weBlue, text width=2.75cm, minimum height=1.0cm] (lerobot) at (0,1.8) {\textbf{LeRobot v3}\\state/action/video};
  \node[wecard=weCyan, text width=2.75cm, minimum height=1.0cm] (rerun) at (0,0.45) {\textbf{Rerun / MCAP}\\multirate logs};
  \node[wecard=weViolet, text width=2.75cm, minimum height=1.0cm] (ncore) at (0,-0.9) {\textbf{NCore}\\sensor calibration};

  \node[wecard=weInk, fill=wePanel, text width=4.05cm, minimum height=3.85cm] (core) at (5.0,0.45) {};
  \node[font=\sffamily\bfseries\large, text=weInk] at (core.north) [yshift=-0.35cm] {WorldEpisode};
  \node[font=\sffamily\small, text=weSlate] at (5.0,1.32) {sidecar blueprint};
  \node[wechip=weBlue] at (3.85,0.65) {identity};
  \node[wechip=weCyan] at (6.05,0.65) {frames + clocks};
  \node[wechip=weAmber] at (3.85,-0.05) {action contract};
  \node[wechip=weGreen] at (6.05,-0.05) {roles + events};
  \node[wechip=weRose] at (4.85,-0.75) {provenance + lineage};

  \node[wecard=weGreen, text width=2.75cm, minimum height=1.0cm] (usd) at (10.0,1.8) {\textbf{OpenUSD}\\world composition};
  \node[wecard=weAmber, text width=2.75cm, minimum height=1.0cm] (gltf) at (10.0,0.45) {\textbf{glTF / SPZ}\\Gaussian splats};
  \node[wecard=weRose, text width=2.75cm, minimum height=1.0cm] (sim) at (10.0,-0.9) {\textbf{MuJoCo / Isaac}\\replay targets};

  \draw[wearrow=weBlue] (lerobot.east) -- ($(core.west)+(0,1.0)$);
  \draw[wearrow=weCyan] (rerun.east) -- (core.west);
  \draw[wearrow=weViolet] (ncore.east) -- ($(core.west)+(0,-1.0)$);
  \draw[wearrow=weGreen] ($(core.east)+(0,1.0)$) -- (usd.west);
  \draw[wearrow=weAmber] (core.east) -- (gltf.west);
  \draw[wearrow=weRose] ($(core.east)+(0,-1.0)$) -- (sim.west);

  \node[wechip=weInk, fill=wePanel, text width=9.6cm, align=center] at (5,-2.05)
    {Executable outputs: validator diagnostics \quad conversion-loss reports \quad lineage-safe splits \quad replay assumptions};
\end{tikzpicture}%
}
\caption{WorldEpisode is the semantic bridge between robot-learning containers, reconstructed
world assets, and simulator targets. It preserves the world--episode contract without replacing
the storage formats on either side.}
\label{fig:worldepisode-bridge}
\end{figure}

This architecture gives the paper its central rule: keep native containers for the data they are
good at storing, but make the cross-container assumptions explicit, checkable, and loss-aware.
